\documentclass[11pt]{article}

\usepackage{amsmath,amssymb,amsthm}
\usepackage{geometry}
\usepackage{hyperref}
\usepackage{enumitem}
\usepackage{mathtools}

\geometry{margin=1in}

\title{Standing Localization and the Structural Inevitability of Life\\
\large --- A Conditional Constraint on Explanatory Physical Regimes}
\author{Amos Jay Maley}
\date{}

\begin{document}
\maketitle

\begin{abstract}
It is often taken for granted that the emergence of life is contingent, accidental, or dependent on highly specific initial conditions. This paper argues for a weaker but more robust claim: under minimal structural constraints required for physical descriptions to retain explanatory standing, life-like structures are unavoidable. The argument does not appeal to teleology, anthropic selection, or biological exceptionalism. Instead, it proceeds from the requirement that viable physical descriptions must license stable reference to localized, persistence-capable structures. We show that, once global coordination is excluded and localization is enforced, persistence under perturbation necessitates internal regulation. Such regulation is structurally equivalent to minimal life. The result is conditional but unavoidable: any physical regime that supports localized standing must admit life-like organization.
\end{abstract}

\section{Introduction}

Why does life appear at all? Traditional answers appeal to chance, selection effects, or special initial conditions. These explanations often presuppose a background physical regime in which localized, persistent structures already exist. This paper takes a different approach. Rather than asking how life arises within a given physical description, we ask what constraints any physically intelligible description must satisfy in order to function as an explanation at all.

The central claim is modest but strong. We do not claim that life is guaranteed, probable, or frequent. We claim only the following: if a physical description is to retain explanatory standing---if it is to successfully refer to stable structures across admissible reformulations---then it must permit localized structures capable of persistence under perturbation. Such persistence, we argue, requires internal regulation. That requirement is already sufficient to force life-like organization.

The argument is structural rather than historical. It does not depend on Earth-specific chemistry, evolutionary pathways, or biological mechanisms. ``Life'' is used minimally, as shorthand for a class of persistence-regulating structures, not as a biological or teleological category.

\paragraph{Standing and terminology (hardening).}
Explanatory standing, as used here, is not a normative ideal imposed on physical description, but a minimal constraint implicit in explanatory practice itself. Descriptions that fail to stabilize reference under admissible reformulation do not merely explain poorly; they cease to function as explanations at all, because they cannot support counterfactual tracking, re-identification, or cross-context comparison. The constraints employed in this paper---localization, exclusion of global coordination, and persistence under perturbation---are therefore not optional methodological preferences, but conditions under which explanatory discourse remains intelligible. Likewise, the term ``life'' is not introduced to elevate biological systems to privileged status, but to name the minimal structural class forced by these constraints. Readers uncomfortable with the term may substitute ``internally regulated persistence-capable structures'' without loss of content. The inevitability claim advanced is thus neither metaphysical nor biological: it is a conditional structural consequence of requiring that physical descriptions retain explanatory standing at all.

\section{Scope and Commitments}

This paper adopts the following commitments and exclusions explicitly:

\begin{itemize}[leftmargin=2em]
\item No teleological assumptions are made.
\item No anthropic selection principles are invoked.
\item No claims are made about the probability, frequency, or typicality of life.
\item ``Life'' is treated as a minimal structural predicate, not a biological taxonomy.
\end{itemize}

The argument relies only on constraints required for physical descriptions to support stable reference and localized persistence. Readers need not accept any broader metaphysical framework.

\section{Operational Definitions}

To keep the argument inspectable, we introduce minimal operational definitions. These are deliberately weaker than many philosophical treatments.

\subsection*{Definition 1 (Standing)}
A physical description has \emph{standing} iff it licenses stable, re-identifiable reference to structures across admissible redescriptions.

\paragraph{Clarification.}
Standing is a constraint on descriptions, not a presupposition about what structures exist. A description may fail to have standing even if structures exist, if reference to them cannot be stabilized under admissible reformulation.

\subsection*{Definition 2 (Localization)}
A structure is \emph{localized} iff its persistence conditions are bounded by finite spatial and energetic constraints.

\subsection*{Definition 3 (Viable Structure)}
A structure is \emph{viable} iff it can maintain identity under environmental perturbation without requiring global coordination.

These definitions are practice-level constraints: they articulate what physical descriptions must enable in order to function explanatorily.

\section{The Constraint of Localization}

Global coordination is incompatible with localization. Any structure whose persistence depends on system-wide synchronization fails to qualify as localized. Physical descriptions that rely on such coordination cannot support stable reference, because perturbations propagate without bound.

Therefore, any structure admitted by a standing-preserving description must be capable of maintaining itself using only locally available resources and interactions.

\paragraph{Lemma (Non-Standing of Globally Coordinated Regimes).}
Any physical regime whose structural persistence depends on system-wide coordination fails to support explanatory standing. In such regimes, reference stability becomes sensitive to unbounded perturbations, preventing localized explanatory claims from remaining invariant under admissible reformulation.

This exclusion is structural rather than pragmatic: global coordination undermines the very conditions under which reference to localized structures can be stabilized.

\section{Persistence Under Perturbation}

Purely inert localized structures are fragile. In realistic environments, perturbations are unavoidable. A structure that merely persists passively will decohere, disperse, or collapse.

Persistence under perturbation therefore requires active stabilization. This stabilization must be internal to the structure, since external global control is excluded by localization.

Thus, viability entails internal regulation.

\section{Internal Regulation and Structural Life}

Internal regulation requires:
\begin{enumerate}[leftmargin=2em]
\item selective exchange with the environment,
\item maintenance of boundary conditions,
\item feedback-sensitive response to perturbations.
\end{enumerate}

These features are not biologically specific. They are structural necessities for any localized system that persists non-trivially.

\paragraph{Lemma (Minimal Life Equivalence).}
Any structure exhibiting selective exchange, boundary maintenance, and perturbation-responsive regulation satisfies the minimal criteria for life employed in this paper. No biological, chemical, or evolutionary assumptions are required beyond these structural features.

Together, these properties constitute a minimal life-like organization: metabolism-like exchange, boundary maintenance, and regulatory feedback.

\section{Conditional Inevitability Theorem}

We can now state the central result explicitly.

\paragraph{Theorem (Conditional Inevitability of Life-Like Structures).}
Any physical regime that:
\begin{enumerate}[leftmargin=2em]
\item supports localized structures,
\item excludes global coordination,
\item preserves explanatory standing via stable reference,
\end{enumerate}
must admit life-like structures.

\paragraph{Proof Sketch.}
Localization excludes global control. Exclusion of global control makes passive persistence impossible. Passive persistence failure forces internal regulation. Internal regulation entails selective exchange, boundary maintenance, and feedback. These jointly define minimal life-like organization.

This argument is structural rather than empirical: it proceeds from constraints on admissible description rather than contingent physical facts. Therefore, any regime satisfying the stated constraints must admit such structures. \hfill $\square$

\section{What This Result Does Not Claim}

For clarity, we emphasize the limits of the argument.

\begin{itemize}[leftmargin=2em]
\item It is not a theory of biological origins.
\item It does not claim life is common or inevitable in a probabilistic sense.
\item It does not depend on Earth-like chemistry.
\item It does not imply intentionality, purpose, or design.
\end{itemize}

The result concerns structural admissibility, not empirical frequency.

\section{Conclusion}

If physical descriptions are to retain standing---if they are to explain anything at all---they must support localized, persistence-capable structures. Persistence under perturbation forces internal regulation. Internal regulation forces life-like organization. This chain is structural, not historical.

Life, in its minimal sense, is not an accident added to physics. It is what localized explanatory coherence looks like.

\end{document}
